\begin{frame}[plain]
\wrsTitlePage{Diffusion feature caching}{What do LearniBridge and LinCa change, and where do we stand?}{Week 38 \textperiodcentered{} two new learned competitors (ICML 2026, ECCV 2026) profiled against our schedules, plus our matched runs at N=5-25.}{Weekly competitor watch \textperiodcentered{} three sources \textperiodcentered{} our own runs are labelled separately}
\note{Frame the meeting: this is a competitor read-out, not a paper explainer. Every claim is tagged as the paper's reported result or as our own run.}
\end{frame}

\begin{frame}{Where we are, and what changed this week}
\begin{columns}[T]
\column{0.5\textwidth}
\wrsSection{Established (unchanged)}
\begin{itemize}
  \item \textbf{P1} Block outputs are temporally coherent, so a cached feature is a usable stand-in.
  \item \textbf{Family A} Fixed reuse (DeepCache, FORA) and calibrated reuse (TeaCache, ToCA).
  \item \textbf{Family B} History extrapolation with a fixed rule (TaylorSeer) at longer cache intervals.
\end{itemize}
\column{0.46\textwidth}
\wrsSection{Where we are now}
\wrsBody{Two 2026 competitors move the correction itself into learned parameters, and both report quality that fixed rules do not reach at 5-7x. We now have matched runs against one of them.}\par\vspace{0.4em}
\wrsTiny{Last week: schedule tuning only, no learned competitor in our comparison.}
\end{columns}
\note{Twenty seconds. The delta is the arrival of learned competitors and the first matched run, not a change in our own method.}
\end{frame}

\begin{frame}{Four families, one schedule: where the competitors sit}
\begin{columns}[T]
\column{0.22\textwidth}
\wrsTiny{A}\par\vspace{0.15em}
{\small\bfseries\color{wrsInk}Fixed reuse}\par\vspace{0.3em}
\wrsSmall{DeepCache, FORA reuse a cached block output directly; TeaCache and ToCA add calibrated or dynamic caching decisions.}
\column{0.22\textwidth}
\wrsTiny{B}\par\vspace{0.15em}
{\small\bfseries\color{wrsInk}History forecast}\par\vspace{0.3em}
\wrsSmall{TaylorSeer extrapolates to the next timestep from several cached features with a fixed Taylor-style rule.}
\column{0.22\textwidth}
\wrsTiny{C}\par\vspace{0.15em}
{\small\bfseries\color{wrsInk}Learned correction}\par\vspace{0.3em}
\wrsSmall{LearniBridge fits a low-rank weight correction on the final block so a cached input reproduces the skipped full computation.}
\column{0.22\textwidth}
\wrsTiny{D}\par\vspace{0.15em}
{\small\bfseries\color{wrsInk}Learned decomposition}\par\vspace{0.3em}
\wrsSmall{LinCa learns an invertible feature decomposition and matches a different prediction order to each component, per model and timestep segment.}
\end{columns}
\vspace{0.4em}
\begin{wrsband}{gap}{Open gap}Families C and D learn the correction; both still assume the caching schedule is given. That is where our schedule work sits.\end{wrsband}
\note{Keep the point structural: the field has moved from hand-designed rules to learned corrections. Both new competitors are complementary to a schedule, not replacements for one.}
\end{frame}

\begin{frame}{Competitor 1: LearniBridge (ICML 2026)}
\begin{columns}[T]
\column{0.31\textwidth}
\wrsSection{A. Intuition}
\wrsSmall{Ask what correction makes the cached computation reproduce the skipped full computation, instead of forecasting the future feature.}
\column{0.31\textwidth}
\wrsSection{B. Mechanism}
\wrsSmall{Fit a low-rank LoRA update on the linear layers of the final Transformer block; at skipped steps only that block runs. Calibration uses 3-5 prompts.}
\column{0.31\textwidth}
\wrsSection{C. Where it fails}
\begin{itemize}
  \item Needs a pre-calibration pass with full computation to record targets.
  \item Only the final block is calibrated, so it absorbs the shift of every bypassed block.
\end{itemize}\vspace{0.3em}
\begin{wrsband}{interpretation}{Relation to us}Its adapters presuppose a cache schedule; it changes which block runs, not when skips happen.\end{wrsband}
\end{columns}
\vfill\wrsCitation{Huang et al., arXiv:2606.26778, ICML 2026}
\note{One sentence for the evidence: SVD of the measured input matrix decays fast, and corrections fitted on disjoint prompt groups have small principal angles. Do not restate the derivation; this is a competitor profile.}
\end{frame}

\begin{frame}{Competitor 2: LinCa (ECCV 2026)}
\begin{columns}[T]
\column{0.31\textwidth}
\wrsSection{A. Intuition}
\wrsSmall{Inside one feature, dimensions do not share the same temporal behaviour, so one prediction order for the whole feature is a mismatch.}
\column{0.31\textwidth}
\wrsSection{B. Mechanism}
\wrsSmall{A lightweight invertible network splits the feature into components with distinct continuity, applies a matched prediction order to each, and maps them back with the exact inverse (under 0.2\% extra parameters).}
\column{0.31\textwidth}
\wrsSection{C. Where it fails}
\begin{itemize}
  \item Needs trained decomposition and predictors per model and timestep segment.
  \item Reported near-lossless 5-7x is on FLUX, Qwen-Image and HunyuanVideo; not reproduced by us.
\end{itemize}\vspace{0.3em}
\begin{wrsband}{interpretation}{Relation to us}Our skip pattern still decides how far its prediction must reach.\end{wrsband}
\end{columns}
\vfill\wrsCitation{Liu et al., arXiv:2608.17973, ECCV 2026}
\note{Emphasise the two distinct kinds of specialisation: within a feature (sub-components, different orders) and across the trajectory (per-segment predictors). These are easy to blur.}
\end{frame}

\begin{frame}{The gap both competitors leave}
\begin{wrsband}{gap}{Gap}Both competitors learn a better correction, but the caching schedule stays a fixed, hand-set interval; neither reports what happens when the schedule is data-dependent.\end{wrsband}
\begin{columns}[T]
\column{0.54\textwidth}
\wrsSection{Evidence}
\begin{itemize}
  \item LearniBridge evaluates a fixed full-compute interval N and calibrates one bridge for it.
  \item LinCa fixes the decomposition once per model and trains per segment, with the interval set externally.
  \item In our matched runs, the ranking between methods changes with the interval rather than staying fixed.
\end{itemize}
\column{0.4\textwidth}
\wrsSection{Implication}
\wrsSmall{A schedule-aware competitor comparison is the right test, and it is the piece we can own.}
\end{columns}
\note{This is the pivot from competitor reading to our own work. Keep it hypothesis-shaped; we have not shown that a data-dependent schedule wins.}
\end{frame}

\begin{frame}{Our approach: retarget the cached state, not the network}
\begin{columns}[T]
\column{0.3\textwidth}
\wrsSection{Forecast (TaylorSeer)}
\wrsSmall{Z\_s $\rightarrow$ Z\_d: predict the future feature first, then hope the frozen tail accepts it.}
\column{0.3\textwidth}
\wrsSection{Learned weights (LearniBridge)}
\wrsSmall{F\_W $\rightarrow$ F\_W+$\Delta$W: adapt the weights of the final block to the stale input.}
\column{0.32\textwidth}
\wrsSection{The idea}
\wrsBody{Ours keeps F\_W frozen and edits the hidden state instead. Find Z\textasciitilde{}\_s$\rightarrow$d such that the frozen tail, given the true target condition, produces the target output.}
\end{columns}
\vspace{0.6em}\wrsBody{The problem is not an incorrect timestep label; it is a stale hidden state. At the skipped step the tail still receives the correct target-step condition, so only the state is wrong.}
\wrsMath{Z_s \neq Z_d, \qquad F_{\text{tail}}(Z_s, c_d) \neq F_{\text{tail}}(Z_d, c_d)}
\wrsMath{F_{\text{tail}}^{\text{frozen}}(\tilde Z_{s \to d}, c_d) \approx F_{\text{tail}}^{\text{frozen}}(Z_d, c_d)}
\note{This is the pivot slide for our own work. Say the mechanism in one line: the tail is already correct, the state it receives is stale, so we repair the state and leave the network alone.}
\end{frame}

\begin{frame}{The retargeter: three inputs, one small MLP}
\begin{columns}[T]
\column{0.5\textwidth}
\wrsSection{Three inputs per token}
\begin{itemize}
  \item \textbf{Local.} the cached token state z\_i at step s.
  \item \textbf{Global.} an image-wide summary c\_s, the mean of the normalized tokens.
  \item \textbf{Transition.} an embedding e\_t of the step pair (t\_s, t\_d) and the gap $\Delta$.
\end{itemize}
\column{0.46\textwidth}
\wrsSection{Fuse, then predict the correction}
\wrsBody{Normalize the token, fuse the three signals through a SiLU layer, and predict the correction with a two-layer MLP. The retargeted token is the cached token plus that correction.}\par\vspace{0.4em}
\end{columns}
\wrsMath{x_i = \mathrm{RMSNorm}(z_i), \qquad c_s = \frac{1}{L} \sum_{i=1}^{L} x_i}
\wrsMath{e_t = E(t_s, t_d, \Delta), \qquad h_i = \mathrm{SiLU}\!\left(W_z x_i + W_c c_s + e_t\right)}
\wrsMath{\Delta z_i = W_2 \, \mathrm{SiLU}(W_1 h_i), \qquad \tilde z_{s,i} = z_{s,i} + \Delta z_i}
\note{One new object per beat: the three inputs, then the fusion, then the correction. The global summary is what makes the correction depend on the whole image rather than on one token.}
\end{frame}

\begin{frame}{What is frozen, and what the loss actually asks for}
\begin{wrsband}{gap}{Gap}Only the retargeter is trained. Everything below the edited state stays frozen, and the tail still runs with the true target-step condition.\end{wrsband}
\begin{columns}[T]
\column{0.54\textwidth}
\wrsSection{Evidence}
\begin{itemize}
  \item The frozen tail consumes the retargeted state: (H\_d, v\_d) = F\_tail(Z\textasciitilde{}\_s, c\_d).
  \item The training loss matches the teacher velocity; no term pulls the retargeted state toward Z\_d.
  \item The edited state only has to be a useful pre-image for the tail, not a copy of the future state.
\end{itemize}
\column{0.4\textwidth}
\wrsSection{Implication}
\wrsSmall{Metrics that measure distance to Z\_d can punish the method while the tail output improves; judge the retargeter by the tail output, not by state distance.}
\end{columns}
\wrsMath{\mathcal{L}_{\text{vel}} = \lVert \hat v_d - v_d^{*} \rVert_2^2, \qquad \text{no } \lVert \tilde Z_{s \to d} - Z_d \rVert^2 \text{ term}}
\note{This is the design commitment. If someone asks why we do not supervise the state directly, the answer is here: the state is not the deliverable, the frozen tail output is.}
\end{frame}

\begin{frame}{What our matched runs actually show}
\begin{columns}[T]
\column{0.62\textwidth}
\begin{wrsband}{observation}{Observation}At matched cache intervals, Global-Local leads the fidelity metrics (PSNR, SSIM, HPSv2.1) at N=5, 7 and 10, while LearniBridge-100 leads CLIP at N=7 and N=10; at N=25 all three collapse together.\end{wrsband}
\begin{wrsband}{interpretation}{Interpretation}The learned bridge competes on semantic alignment rather than pixel fidelity in our setting, and no method survives the 25-step interval. Our retargeter is trained only through the tail output, so its gain shows up on quality metrics rather than on distance-to-original fidelity.\end{wrsband}
\column{0.34\textwidth}
\wrsSection{Claim}
\wrsClaimId{C1}\quad\wrsStatus{weakened}
\wrsSmall{A learned low-rank calibration dominates our baseline at matched intervals. Our runs do not support this as stated.}
\end{columns}
\note{Say it plainly: the competitor's released configuration is competitive, not dominant, in our matched setting. Report the weak metric too; the audience is the group that will decide whether to keep chasing it.}
\end{frame}

\begin{frame}{Our matched runs: N=5 and N=7}
\begin{center}\scriptsize
\setlength{\tabcolsep}{3pt}
\renewcommand{\arraystretch}{1.15}
\begin{tabular}{lccccccc}
\toprule
Method & \textbf{CLIP $\uparrow$} & \textbf{IR $\uparrow$} & \textbf{HPSv2.1 $\uparrow$} & \textbf{PSNR $\uparrow$} & \textbf{SSIM $\uparrow$} & \textbf{s/ảnh $\downarrow$} & \textbf{Speed $\uparrow$} \\
\midrule
\textcolor{wrsMuted}{Original FLUX} & 32.65 & 1.00 & 30.40 & $\infty$ & 1 & 22.88 & 1.00x \\
Cache thường (N=5) & 32.63 & 0.94 & 29.44 & 14.70 & 0.55 & \textbf{\textcolor{wrsGreen}{4.95}} & 4.63x \\
Global-Local (N=5)\newline{\tiny\color{wrsMuted}(ours)} & \textbf{\textcolor{wrsGreen}{32.79}} & \textbf{\textcolor{wrsGreen}{1.01}} & \textbf{\textcolor{wrsGreen}{30.00}} & \textbf{\textcolor{wrsGreen}{15.21}} & \textbf{\textcolor{wrsGreen}{0.57}} & 4.98 & 4.60x \\
LearniBridge-100 (N=5) & 32.66 & 0.98 & 29.53 & 15.07 & 0.55 & 5.09 & 4.49x \\
Cache thường (N=7) & 32.28 & 0.83 & 28.24 & 14.14 & 0.52 & \textbf{\textcolor{wrsGreen}{4.04}} & 5.66x \\
Global-Local (N=7)\newline{\tiny\color{wrsMuted}(ours)} & 32.63 & 0.92 & \textbf{\textcolor{wrsGreen}{29.10}} & \textbf{\textcolor{wrsGreen}{14.73}} & \textbf{\textcolor{wrsGreen}{0.55}} & 4.09 & 5.60x \\
LearniBridge-100 (N=7) & \textbf{\textcolor{wrsGreen}{32.76}} & \textbf{\textcolor{wrsGreen}{0.92}} & 28.44 & 14.49 & 0.54 & 4.20 & 5.45x \\
\bottomrule
\end{tabular}
\end{center}
\vspace{0.25em}\wrsQuiet{Ours = the retargeter (Global-Local). Green marks the best value in each interval, computed separately for N=5 and N=7; arrows say whether higher or lower is better. PSNR/SSIM are measured against Original FLUX, which is a reference row and is never highlighted.}
\note{The comparison is matched, so we can read across a row group. Do not polish the deltas; several are inside noise and the group should see that.}
\end{frame}

\begin{frame}{Our matched runs: N=10 and N=25}
\begin{center}\scriptsize
\setlength{\tabcolsep}{3pt}
\renewcommand{\arraystretch}{1.15}
\begin{tabular}{lccccccc}
\toprule
Method & \textbf{CLIP $\uparrow$} & \textbf{IR $\uparrow$} & \textbf{HPSv2.1 $\uparrow$} & \textbf{PSNR $\uparrow$} & \textbf{SSIM $\uparrow$} & \textbf{s/ảnh $\downarrow$} & \textbf{Speed $\uparrow$} \\
\midrule
Cache thường (N=10) & 31.72 & 0.66 & 26.36 & 13.77 & 0.48 & \textbf{\textcolor{wrsGreen}{2.70}} & 8.48x \\
Global-Local (N=10)\newline{\tiny\color{wrsMuted}(ours)} & 31.77 & \textbf{\textcolor{wrsGreen}{0.76}} & \textbf{\textcolor{wrsGreen}{27.76}} & \textbf{\textcolor{wrsGreen}{14.46}} & \textbf{\textcolor{wrsGreen}{0.52}} & 2.75 & 8.32x \\
LearniBridge-100 (N=10) & \textbf{\textcolor{wrsGreen}{32.16}} & 0.73 & 26.68 & 14.10 & 0.47 & 2.88 & 7.96x \\
Cache thường (N=25) & 23.08 & -1.48 & 13.71 & 13.12 & 0.42 & \textbf{\textcolor{wrsGreen}{1.36}} & 16.87x \\
Global-Local (N=25)\newline{\tiny\color{wrsMuted}(ours)} & 22.99 & \textbf{\textcolor{wrsGreen}{-1.28}} & \textbf{\textcolor{wrsGreen}{15.68}} & \textbf{\textcolor{wrsGreen}{14.26}} & \textbf{\textcolor{wrsGreen}{0.49}} & 1.41 & 16.24x \\
LearniBridge-100 (N=25) & \textbf{\textcolor{wrsGreen}{23.09}} & -1.51 & 13.72 & 13.47 & 0.45 & 1.55 & 14.77x \\
\bottomrule
\end{tabular}
\end{center}
\vspace{0.25em}\wrsQuiet{Green marks the best value within each interval (N=10, then N=25). At N=25 every method is far from Original FLUX (IR below zero, HPSv2.1 near half), so N=25 is a failure regime for all three and should not be ranked.}
\note{Say explicitly that N=25 is not a ranking, it is a boundary. If anyone quotes a winner there, correct it live.}
\end{frame}

\begin{frame}{Qualitative, N=5: Original / Cache / ours / LearniBridge-100}
\wrsFigure[0.99\linewidth]{assets/N5_qualitative.jpg}
\vspace{0.3em}\wrsQuiet{Sheet split in two blocks: rows 1-4 on the left, rows 5-8 on the right. Columns are Original, Cache (N=5), Global-Local (ours), LearniBridge-100.}
\note{Source: /Users/nguyenthanhlam/LESA/artifacts/learnibridge\_100\_matched/qualitative\_by\_n/N5.jpg, re-composed as two side-by-side halves for slide readability.}
\end{frame}

\begin{frame}{Qualitative, N=7: Original / Cache / ours / LearniBridge-100}
\wrsFigure[0.99\linewidth]{assets/N7_qualitative.jpg}
\vspace{0.3em}\wrsQuiet{Rows 1-4 on the left, rows 5-8 on the right. At N=7 the plain cache starts to drift on structure while both learned methods keep the layout.}
\note{Source: /Users/nguyenthanhlam/LESA/artifacts/learnibridge\_100\_matched/qualitative\_by\_n/N7.jpg. Look for structure drift and colour shifts, not for a winner.}
\end{frame}

\begin{frame}{Qualitative, N=10: Original / Cache / ours / LearniBridge-100}
\wrsFigure[0.99\linewidth]{assets/N10_qualitative.jpg}
\vspace{0.3em}\wrsQuiet{Rows 1-4 on the left, rows 5-8 on the right. Failures are no longer uniform: which method drifts depends on the prompt.}
\note{Source: /Users/nguyenthanhlam/LESA/artifacts/learnibridge\_100\_matched/qualitative\_by\_n/N10.jpg. The prompt dependence is the argument for schedule awareness.}
\end{frame}

\begin{frame}{Qualitative, N=25: Original / Cache / ours / LearniBridge-100}
\wrsFigure[0.99\linewidth]{assets/N25_qualitative.jpg}
\vspace{0.3em}\wrsQuiet{Rows 1-4 on the left, rows 5-8 on the right. At N=25 all four columns are visibly different from Original; this sheet documents the failure regime, not a comparison.}
\note{Source: /Users/nguyenthanhlam/LESA/artifacts/learnibridge\_100\_matched/qualitative\_by\_n/N25.jpg. Use this to set the boundary of the useful interval, then move on.}
\end{frame}

\begin{frame}{What we still do not know}
\begin{columns}[T]
\column{0.5\textwidth}
\wrsSection{Current limitations}
\begin{itemize}
  \item LinCa is profiled from its paper only; we have not reproduced it in our matched setting.
  \item Our runs cover one model (FLUX) and one prompt set, so cross-model ranking is unmeasured.
  \item PSNR/SSIM against Original FLUX penalise every accelerator; they are fidelity, not quality.
  \item N=25 is a failure regime for all three methods, so the interval range where learned corrections pay off is still narrow.
\end{itemize}
\column{0.46\textwidth}
\wrsSection{Open questions}
\begin{itemize}
  \item Does LinCa's learned decomposition close the N=10-25 gap that no current method survives?
  \item Is the CLIP advantage of LearniBridge-100 at N=7/10 reproducible across prompt sets?
  \item Would a data-dependent schedule change which of these competitors wins?
\end{itemize}
\end{columns}
\note{End here. The next action is a LinCa reproduction in the same matched harness; no thank-you slide.}
\end{frame}
